\section{Đối tượng của lí thuyết tập hợp}

\ 

Từ cái tên, chúng ta cũng đã biết lí thuyết tập hợp sẽ tìm hiểu về tính chất của các \defText{tập hợp} (hay \defText{tập}). Chúng ta không định nghĩa tập hợp mà mặc nhiên cho nó tự tồn tại. Thêm vào đó, chúng ta cấp vị từ $\defMath{\in}\hspace{-0.25em}()$ nhận hạng thức là hai tập hợp $x$ và $y$ bất kì. Điều này có nghĩa $\defMath{\in}\hspace{-0.25em}\defMath{(x; y)}$ hay $\defMath{x \in y}$ là một mệnh đề có thể đúng hoặc sai. Nếu sai, chúng ta kí hiệu $\defMath{x \notin y}$. Nếu $x \in y$, chúng ta viết bằng lời ``$x$ thuộc $y$'' hay ``$x$ là một \defText{phần tử} của $y$''. Ngược lại, chúng ta viết ``$x$ không thuộc $y$''. 

Nếu như ghét người đọc, bạn đọc có thể viết theo thứ tự ngược lại:
\begin{align*}
    \defMath{y \ni x} &\iff \defMath{x \in y}; \\
    \defMath{y \not}\hspace{0.08em}\defMath{\ni x} &\iff \defMath{x \notin y}.
\end{align*}
Theo thứ tự, hai mệnh đề được phát biểu là ``$y$ chứa $x$'' và ``$y$ không chứa $x$''.